% ============================================================
%  TD6 -- Benchmarking de Modelos de Deteccion de Fallos por Motor
%  Curso: Machine Learning Industrial
%  Version ampliada con analisis del "por que" de cada decision
% ============================================================
\documentclass[12pt,a4paper]{article}

% ---------- Paquetes ----------
\usepackage[a4paper, top=2.5cm, bottom=2.5cm, left=2.5cm, right=2.5cm]{geometry}
\usepackage[spanish]{babel}
\usepackage[utf8]{inputenc}
\usepackage[T1]{fontenc}
\usepackage{amsmath,amssymb}
\usepackage{booktabs}
\usepackage{graphicx}
\usepackage{xcolor}
\usepackage{listings}
\usepackage{hyperref}
\usepackage{parskip}
\usepackage{titlesec}
\usepackage{enumitem}
\usepackage{array}
\usepackage{longtable}
\usepackage{microtype}
\usepackage{fancyhdr}
\usepackage{lmodern}
\usepackage{tcolorbox}
\usepackage{tabularx}
\usepackage{multirow}
\usepackage{colortbl}
\usepackage{mdframed}

% ---------- Colores ----------
\definecolor{codeblue}{RGB}{0,70,130}
\definecolor{codegray}{RGB}{128,128,128}
\definecolor{codegreen}{RGB}{0,120,0}
\definecolor{codered}{RGB}{160,0,0}
\definecolor{codebg}{RGB}{248,248,248}
\definecolor{titlebg}{RGB}{0,50,100}
\definecolor{whybg}{RGB}{230,245,255}
\definecolor{whyborder}{RGB}{0,100,180}
\definecolor{warnbg}{RGB}{255,245,220}
\definecolor{warnborder}{RGB}{200,140,0}
\definecolor{keybg}{RGB}{235,250,235}
\definecolor{keyborder}{RGB}{0,130,0}
\definecolor{rowgray}{RGB}{242,242,242}

% ---------- Cajas de "Por que" ----------
\newmdenv[
  backgroundcolor=whybg,
  linecolor=whyborder,
  linewidth=2pt,
  leftmargin=0pt,
  rightmargin=0pt,
  innerleftmargin=10pt,
  innerrightmargin=10pt,
  innertopmargin=8pt,
  innerbottommargin=8pt,
  skipabove=8pt,
  skipbelow=8pt
]{porquebox}

\newmdenv[
  backgroundcolor=warnbg,
  linecolor=warnborder,
  linewidth=2pt,
  leftmargin=0pt,
  rightmargin=0pt,
  innerleftmargin=10pt,
  innerrightmargin=10pt,
  innertopmargin=8pt,
  innerbottommargin=8pt,
  skipabove=8pt,
  skipbelow=8pt
]{advertenciabox}

\newmdenv[
  backgroundcolor=keybg,
  linecolor=keyborder,
  linewidth=2pt,
  leftmargin=0pt,
  rightmargin=0pt,
  innerleftmargin=10pt,
  innerrightmargin=10pt,
  innertopmargin=8pt,
  innerbottommargin=8pt,
  skipabove=8pt,
  skipbelow=8pt
]{claveBox}

% ---------- Estilo listings ----------
\lstset{
  backgroundcolor=\color{codebg},
  basicstyle=\ttfamily\small,
  keywordstyle=\color{codeblue}\bfseries,
  commentstyle=\color{codegreen}\itshape,
  stringstyle=\color{codered},
  numberstyle=\tiny\color{codegray},
  breaklines=true,
  breakatwhitespace=true,
  numbers=left,
  numbersep=8pt,
  frame=single,
  framesep=4pt,
  rulecolor=\color{codegray},
  showstringspaces=false,
  tabsize=4,
  captionpos=b,
  language=Python,
  literate={\'a}{{\'{a}}}1 {\'e}{{\'{e}}}1 {\'i}{{\'{i}}}1 {\'o}{{\'{o}}}1 {\'u}{{\'{u}}}1
           {\'A}{{\'{A}}}1 {\'E}{{\'{E}}}1 {\'I}{{\'{I}}}1 {\'O}{{\'{O}}}1 {\'U}{{\'{U}}}1
           {\~n}{{\~{n}}}1 {\~N}{{\~{N}}}1
}

% ---------- Hiperref ----------
\hypersetup{
  colorlinks=true,
  linkcolor=codeblue,
  urlcolor=codeblue,
  citecolor=codeblue,
  pdftitle={TD6: Benchmarking de Modelos de Deteccion de Fallos},
  pdfauthor={Machine Learning Industrial}
}

% ---------- Cabecera/Pie ----------
\pagestyle{fancy}
\fancyhf{}
\fancyhead[L]{\small\textit{TD6 -- Benchmarking de Modelos de Detecci\'{o}n de Fallos}}
\fancyhead[R]{\small\textit{Machine Learning Industrial}}
\fancyfoot[C]{\thepage}
\renewcommand{\headrulewidth}{0.4pt}

% ---------- Formato de secciones ----------
\titleformat{\section}[block]
  {\large\bfseries\color{titlebg}}
  {\thesection.}{0.5em}{}[\titlerule]

\titleformat{\subsection}[block]
  {\normalsize\bfseries\color{codeblue}}
  {\thesubsection.}{0.5em}{}

\titleformat{\subsubsection}[block]
  {\normalsize\itshape\bfseries}
  {\thesubsubsection.}{0.5em}{}

% ============================================================
\begin{document}
% ============================================================

% ---------- PORTADA ----------
\begin{titlepage}
\centering
\vspace*{1.5cm}

{\Large\bfseries Machine Learning Industrial\par}
\vspace{0.3cm}
{\large Data Challenge -- Detecci\'{o}n de Fallos en Motores Rob\'{o}ticos\par}

\vspace{1.2cm}
\hrule height 2pt
\vspace{0.5cm}

{\Huge\bfseries\color{titlebg}
TD6: Benchmarking de Modelos\\[0.3em]
de Detecci\'{o}n de Fallos\\[0.3em]
por Motor
\par}

\vspace{0.5cm}
\hrule height 2pt
\vspace{1.5cm}

\begin{center}
\begin{tabular}{ll}
\textbf{Asignatura:} & Machine Learning Industrial \\[4pt]
\textbf{Trabajo:}    & TD6 -- Validaci\'{o}n Cruzada y Benchmarking \\[4pt]
\textbf{Plataforma:} & Kaggle Data Challenge \\[4pt]
\textbf{Dataset:}    & 6 Motores Servo-Rob\'{o}ticos \\[4pt]
\textbf{Fecha:}      & Junio 2026 \\
\end{tabular}
\end{center}

\vspace{1.5cm}

\begin{abstract}
\noindent
Este documento presenta el benchmarking completo de cuatro modelos de machine learning
para la detecci\'{o}n de fallos individuales en seis motores servo-rob\'{o}ticos, en el marco de
un Data Challenge de Kaggle. Se aplica validaci\'{o}n cruzada de 5 pliegues con b\'{u}squeda
anidada de hiperpar\'{a}metros para comparar AAKR (\textit{Auto-Associative Kernel Regression}),
Random Forest, Gradient Boosting e Histogramado y Regresi\'{o}n Log\'{i}stica, evaluando
precisi\'{o}n, recall y F1-score. Cada decisi\'{o}n t\'{e}cnica --desde el preprocesamiento hasta
la arquitectura de evaluaci\'{o}n-- se analiza en profundidad explicando el \textit{por qu\'{e}}
de cada elecci\'{o}n. Se selecciona el mejor modelo por motor, se entrena sobre el conjunto
completo de entrenamiento y se genera una \textit{submission} para Kaggle. El an\'{a}lisis
revela que Random Forest con pesos balanceados domina gracias a su robustez al desbalanceo
extremo de clases (hasta 200:1), y explica mediante an\'{a}lisis estad\'{i}stico la raz\'{o}n
por la que la alta varianza entre pliegues es un fen\'{o}meno esperado y no un defecto
del dise\~{n}o experimental.
\end{abstract}

\vfill
{\small Generado con \LaTeX{}}
\end{titlepage}

% ---------- INDICE ----------
\tableofcontents
\newpage

% ============================================================
\section{Introducci\'{o}n: Por qu\'{e} este TD}
% ============================================================

\subsection{Contexto y Motivaci\'{o}n del Problema}

El TD6 representa un salto cualitativo respecto a los trabajos anteriores del curso:
se pasa de experimentos te\'{o}ricos sobre datasets sint\'{e}ticos (TD5) a un \textbf{problema
real de ingenier\'{i}a industrial} con todas sus complejidades. Este salto merece ser
comprendido en detalle.

\begin{porquebox}
\textbf{Por qu\'{e} este TD es distinto a los anteriores}\\[4pt]
En TD5 se trabaj\'{o} con distribuciones controladas y balanceadas. El TD6 enfrenta
problemas que aparecen en sistemas industriales reales:
\begin{itemize}[itemsep=2pt]
  \item \textbf{6 motores con tasas de fallo heterog\'{e}neas}: desde 0.4\% hasta 17\%.
    No existe un modelo universal; cada motor puede necesitar una estrategia distinta.
  \item \textbf{Se\~{n}ales no estacionarias}: las condiciones de operaci\'{o}n cambian entre
    experimentos (temperatura ambiente diferente, calibraciones distintas).
  \item \textbf{Objetivo de generalizaci\'{o}n real}: el modelo se evaluar\'{a} en datos de test
    nunca vistos durante el desarrollo, en condiciones que pueden diferir del
    entrenamiento.
  \item \textbf{Restricci\'{o}n de recursos}: el tiempo de c\'{o}mputo de Kaggle limita
    la profundidad de b\'{u}squeda de hiperpar\'{a}metros.
\end{itemize}
\end{porquebox}

\subsection{Estructura del Dataset y sus Implicaciones}

El dataset proviene de un sistema de manipulaci\'{o}n industrial con seis motores
servo-rob\'{o}ticos, cada uno equipado con tres sensores:

\begin{center}
\begin{tabular}{llll}
\toprule
\textbf{Sensor} & \textbf{Unidad} & \textbf{Rango f\'{i}sico} & \textbf{Significado f\'{i}sico} \\
\midrule
Posici\'{o}n & Conteos encoder & $[0,\, 1000]$ & \'{A}ngulo del eje motor \\
Temperatura & $^\circ$C & $[0,\, 150]$ & Calor generado por fricci\'{o}n/corriente \\
Voltaje & mV & $[6\,000,\, 9\,000]$ & Tensi\'{o}n de alimentaci\'{o}n del driver \\
\bottomrule
\end{tabular}
\end{center}

Las dimensiones del dataset son:

\begin{center}
\begin{tabular}{lrrl}
\toprule
\textbf{Partici\'{o}n} & \textbf{Filas} & \textbf{Secuencias} & \textbf{Frecuencia de muestreo} \\
\midrule
Entrenamiento & 39\,309 & 23 & 10 Hz (0.1 s/muestra) \\
Test          & 14\,157 & 8  & 10 Hz (0.1 s/muestra) \\
\bottomrule
\end{tabular}
\end{center}

\begin{advertenciabox}
\textbf{Implicaci\'{o}n cr\'{i}tica del dise\~{n}o del dataset}\\[4pt]
Con solo 23 secuencias en entrenamiento y 8 en test, el dataset tiene muy poca
diversidad a nivel de \textit{experimento completo}. Esto significa que:
\begin{itemize}[itemsep=2pt]
  \item La validaci\'{o}n cruzada de 5 pliegues opera con $\approx 4{-}5$ secuencias
    en cada pliegue de validaci\'{o}n.
  \item Si los fallos se concentran en 1--2 secuencias, la aparici\'{o}n o no de esas
    secuencias en el pliegue de validaci\'{o}n domina el resultado del fold.
  \item Esto explica la alta varianza entre pliegues observada en los resultados.
\end{itemize}
Esta no es una limitaci\'{o}n del m\'{e}todo elegido; es una propiedad del dataset.
\end{advertenciabox}

\subsection{Distribuci\'{o}n de Fallos por Motor}

La siguiente tabla resume las tasas de fallo aproximadas seg\'{u}n los resultados
observados en entrenamiento:

\begin{center}
\begin{tabular}{cccll}
\toprule
\textbf{Motor} & \textbf{Fallos (\%)} & \textbf{Ratio desbalanceo} & \textbf{Dificultad} & \textbf{Implicaci\'{o}n} \\
\midrule
Motor 1 & $\sim 4{-}5\%$  & $\sim 20:1$  & Media    & CV estable, F1 alcanzable \\
Motor 2 & $\sim 17\%$     & $\sim 5:1$   & Baja     & Dataset m\'{a}s balanceado \\
Motor 3 & $< 0.5\%$       & $>200:1$     & Extrema  & Clasificadores colapsan \\
Motor 4 & $\sim 17\%$     & $\sim 5:1$   & Media    & Error de ejecuci\'{o}n en notebook \\
Motor 5 & Desconocida     & Desconocido  & --       & Error de ejecuci\'{o}n en notebook \\
Motor 6 & Desconocida     & Desconocido  & --       & Error de ejecuci\'{o}n en notebook \\
\bottomrule
\end{tabular}
\end{center}

\subsection{Objetivos del TD6}

\begin{enumerate}[itemsep=3pt]
  \item Aplicar \textbf{validaci\'{o}n cruzada k-fold anidada} ($k_{\text{ext}}=5$,
    $k_{\text{int}}=3$) para comparar cuatro modelos de detecci\'{o}n de fallos en
    cada uno de los seis motores.
  \item Seleccionar el \textbf{mejor modelo} por motor seg\'{u}n el F1-score macro.
  \item Entrenar el modelo seleccionado sobre el conjunto completo de entrenamiento y
    generar una \textbf{predicci\'{o}n} para el conjunto de test.
  \item Realizar la \textbf{submission a Kaggle} en el formato requerido.
  \item Analizar las \textbf{causas de varianza} en los resultados de validaci\'{o}n cruzada
    y justificar cada decisi\'{o}n de dise\~{n}o del pipeline.
\end{enumerate}

% ============================================================
\section{Preprocesamiento: Justificaci\'{o}n Profunda de Cada Paso}
% ============================================================

\subsection{Visi\'{o}n General del Pipeline de Preprocesamiento}

Antes de entrenar cualquier modelo, los datos crudos pasan por una funci\'{o}n unificada
de preprocesamiento. El orden de las operaciones es intencional:

\begin{lstlisting}[caption={Funci\'{o}n principal de preprocesamiento}, label={lst:preproc}]
def pre_processing(df):
    remove_outliers(df)        # Paso 1: Recorte por rango fisico + forward-fill
    compensate_seq_bias(df)    # Paso 2: Restar valor inicial de cada secuencia
    cal_diff(df, 10)           # Paso 3: Diferencias con lag 10
    df.fillna(0, inplace=True) # Paso 4: Rellenar NaN residuales con 0
\end{lstlisting}

El diagrama conceptual del flujo de transformaci\'{o}n de una se\~{n}al es:

\begin{center}
\begin{verbatim}
Señal cruda:     [35.1, 35.3, 35.8, 36.5, 120.0, 36.2, 36.4, ...]
                                              ^
                                       Outlier fisico
                                       (sensor saturado)
          |
          | remove_outliers()
          v
Clip+Fill:       [35.1, 35.3, 35.8, 36.5, 36.5,  36.2, 36.4, ...]
                 (120.0 -> NaN -> forward-fill con 36.5)
          |
          | compensate_seq_bias()
          v
Relativa:        [ 0.0,  0.2,  0.7,  1.4,  1.4,   1.1,  1.3, ...]
                 (restado x[0]=35.1; la señal ahora mide desviacion)
          |
          | cal_diff(lag=10)
          v
Dinamica:        [ NaN, NaN, ..., 0.2+..., ...]
                 (diferencia a 10 pasos: captura velocidad de cambio)
          |
          | fillna(0)
          v
Final:           [ 0.0, 0.0, ..., 0.8, 1.2, ...]
                 (NaN iniciales en diferencias = 0 = sin cambio detectado)
\end{verbatim}
\end{center}

\subsection{Paso 1: Eliminaci\'{o}n de Outliers}

\subsubsection{Descripci\'{o}n T\'{e}cnica}

Los sensores industriales producen ocasionalmente lecturas f\'{i}sicamente imposibles
por errores de transmisi\'{o}n, saturaci\'{o}n del ADC o falta de actualizaci\'{o}n del bus.
La funci\'{o}n aplica \textit{clipping} por rango f\'{i}sico conocido:
\begin{itemize}[itemsep=2pt]
  \item Posici\'{o}n: valores fuera de $[0, 1000]$ conteos se marcan como \texttt{NaN}.
  \item Temperatura: valores negativos o superiores a 150\textdegree{}C se eliminan.
  \item Voltaje: valores fuera de $[6\,000, 9\,000]$ mV se eliminan.
\end{itemize}
Los \texttt{NaN} resultantes se rellenan con \textit{forward-fill} (propagaci\'{o}n del
\'{u}ltimo valor v\'{a}lido), preservando la continuidad temporal de la se\~{n}al.

\begin{porquebox}
\textbf{Por qu\'{e} forward-fill en lugar de interpolaci\'{o}n lineal o eliminar filas}\\[4pt]
\textbf{Opci\'{o}n 1: Eliminar filas con outliers.} Introduce huecos en la serie
temporal, rompiendo la continuidad. Los modelos basados en diferencias (diff) y en
ventanas deslizantes asumen continuidad; un hueco produce diferencias artefactuales
enormes ($\approx$ 10--50x el valor real), contaminando las features din\'{a}micas.\\[4pt]
\textbf{Opci\'{o}n 2: Interpolaci\'{o}n lineal.} Requiere conocer los valores futuros
(acausalidad), lo que no es posible en sistemas en tiempo real. Para ser causal,
habr\'{i}a que hacer backward-fill, que tiene el mismo efecto pr\'{a}ctico que forward-fill.\\[4pt]
\textbf{Opci\'{o}n 3: Forward-fill (elegida).} Mantiene el \'{u}ltimo valor v\'{a}lido,
lo que es equivalente a asumir ``el sensor est\'{a} congelado'' durante el error de
transmisi\'{o}n. Preserva la continuidad temporal y es causal.
\end{porquebox}

\subsection{Paso 2: Compensaci\'{o}n de Sesgo por Secuencia}

\subsubsection{Descripci\'{o}n T\'{e}cnica}

Distintos experimentos pueden comenzar con condiciones iniciales diferentes. Por
ejemplo, la temperatura inicial puede variar seg\'{u}n la temperatura ambiente del
laboratorio ese d\'{i}a, o la posici\'{o}n inicial puede diferir seg\'{u}n la \'{u}ltima
tarea ejecutada. La transformaci\'{o}n es:

\begin{equation}
  x'_t = x_t - x_0, \quad \forall t \in \text{secuencia}
\end{equation}

donde $x_0$ es el valor de la se\~{n}al en el primer instante de esa secuencia.

\begin{porquebox}
\textbf{Por qu\'{e} restar el valor inicial (compensar sesgo de secuencia)}\\[4pt]
\textbf{Problema sin compensaci\'{o}n}: Supongamos que el Motor 1 en el experimento
\texttt{20240101} empieza a 25\textdegree{}C y en el experimento \texttt{20240325}
empieza a 35\textdegree{}C. Un fallo en ambos experimentos ocurre cuando la temperatura
sube 15\textdegree{}C respecto al inicio (patr\'{o}n relativo). Sin compensaci\'{o}n, el
modelo aprender\'{i}a:
\begin{itemize}[itemsep=2pt]
  \item ``Temperatura 40\textdegree{}C = fallo'' (para el experimento de enero)
  \item ``Temperatura 50\textdegree{}C = fallo'' (para el experimento de marzo)
\end{itemize}
Esto es inconsistente: el mismo tipo de fallo aparece en valores absolutos distintos.
El modelo necesita muchos m\'{a}s datos para generalizar.\\[4pt]
\textbf{Con compensaci\'{o}n}: En ambos experimentos, el fallo ocurre cuando
$x'_t = +15^\circ\text{C}$. El patr\'{o}n es ahora \textbf{invariante} a las
condiciones iniciales del experimento. El modelo puede generalizar con mucho
menos datos.\\[4pt]
\textbf{Analog\'{i}a}: Es como medir la velocidad de un objeto en lugar de su posici\'{o}n
absoluta. Dos coches pueden colisionar aunque partan de posiciones distintas; lo que
importa es la tasa de cambio, no el valor absoluto.
\end{porquebox}

\subsubsection{Efecto Cuantitativo Esperado}

\begin{center}
\begin{tabular}{lcc}
\toprule
\textbf{Escenario} & \textbf{Sin compensaci\'{o}n} & \textbf{Con compensaci\'{o}n} \\
\midrule
Fallo en experimento con $T_0=25^\circ$C & $T=40^\circ$C & $\Delta T=+15^\circ$C \\
Fallo en experimento con $T_0=35^\circ$C & $T=50^\circ$C & $\Delta T=+15^\circ$C \\
Patr\'{o}n que aprende el modelo & 2 patrones distintos & 1 patr\'{o}n \'{u}nico \\
Muestras necesarias para generalizar & $O(N^2)$ & $O(N)$ \\
\bottomrule
\end{tabular}
\end{center}

\subsection{Paso 3: Diferencias con Lag 10}

\subsubsection{Descripci\'{o}n T\'{e}cnica}

La diferencia con lag $\ell = 10$ captura la \textit{tasa de cambio} de cada se\~{n}al
a la escala de 10 pasos temporales (= 1 segundo a 10 Hz):

\begin{equation}
  \Delta x_t^{(\ell)} = x_t - x_{t-\ell}
\end{equation}

Se calculan para temperatura, voltaje y posici\'{o}n de cada motor, a\~{n}adiendo
$3 \times 6 = 18$ columnas nuevas al dataframe.

\begin{porquebox}
\textbf{Por qu\'{e} diferencias en lugar de usar valores absolutos solamente}\\[4pt]
\textbf{Informaci\'{o}n de los valores absolutos}: captura el \textit{estado actual}
del motor. Un voltaje de 7500 mV es normal; 6100 mV puede indicar un problema el\'{e}ctrico.\\[4pt]
\textbf{Informaci\'{o}n de las diferencias}: captura la \textit{din\'{a}mica} del motor.
Una temperatura que sube 5\textdegree{}C en 1 segundo indica sobrecalentamiento
inminente aunque el valor absoluto actual sea ``normal''. Una ca\'{i}da brusca de
voltaje en 0.5 segundos indica un cortocircuito transitorio.\\[4pt]
\textbf{Sin diferencias}: el modelo solo ve valores absolutos y puede no distinguir
entre un motor que lleva horas a 45\textdegree{}C (estable, posiblemente normal) y
uno que acaba de llegar a 45\textdegree{}C desde 30\textdegree{}C en 2 segundos
(en r\'{a}pido calentamiento, posiblemente fallando).\\[4pt]
\textbf{Con diferencias}: el modelo ve ambos aspectos: el estado ($x_t$) y la
velocidad de cambio ($\Delta x_t^{(10)}$). Esto duplica la informaci\'{o}n
diagn\'{o}stica disponible.
\end{porquebox}

\subsubsection{Justificaci\'{o}n de la Elecci\'{o}n de Lag = 10}

La elecci\'{o}n del lag es un compromiso entre tres factores:

\begin{center}
\begin{tabular}{lccc}
\toprule
\textbf{Lag} & \textbf{Escala temporal} & \textbf{Sensibilidad a ruido} & \textbf{Sensibilidad a cambios r\'{a}pidos} \\
\midrule
$\ell = 1$  & 0.1 s & Muy alta (ruido dominante) & Muy alta (puede ser ruido) \\
$\ell = 5$  & 0.5 s & Alta & Alta \\
$\ell = 10$ & 1.0 s & Media (compromiso) & Media-alta \\
$\ell = 50$ & 5.0 s & Baja & Baja (pierde eventos r\'{a}pidos) \\
$\ell = 100$& 10.0 s & Muy baja & Muy baja \\
\bottomrule
\end{tabular}
\end{center}

\begin{porquebox}
\textbf{Por qu\'{e} lag = 10 (1 segundo)} \\[4pt]
1 segundo es una escala temporal naturalmente \'{u}til para detectar fallos en motores
servo-rob\'{o}ticos porque:
\begin{itemize}[itemsep=2pt]
  \item Los fallos el\'{e}ctricos (cortocircuitos, sobreintensidades) se manifiestan en
    escalas de 0.1--1 segundo: lag=1 o lag=10.
  \item Los fallos mec\'{a}nicos (sobrecalentamiento, vibraci\'{o}n excesiva) se manifiestan
    en escalas de 1--10 segundos: lag=10 o lag=100.
  \item Lag=10 es el punto medio que captura ambas categor\'{i}as.
  \item En 1 segundo, la temperatura no puede subir significativamente por ruido del
    sensor (ruido t\'{i}pico: $\pm 0.1^\circ$C); cualquier cambio es real.
\end{itemize}
\end{porquebox}

\subsection{Features Utilizadas por Tipo de Modelo}

\begin{center}
\begin{tabularx}{\textwidth}{lXX}
\toprule
\textbf{Tipo de modelo} & \textbf{Features (total)} & \textbf{Justificaci\'{o}n} \\
\midrule
Clasificadores (RF, GB, LR) &
  Temperatura $\times 6$ motores + Voltaje $\times 6$ motores + Diffs posici\'{o}n, temp., voltaje $\times 6$ motores = \textbf{30 features} &
  No incluye posici\'{o}n directa porque la posici\'{o}n absoluta no indica fallo; su din\'{a}mica (diff) s\'{i} \\
\addlinespace
AAKR (Regresor) &
  \texttt{time} + Posici\'{o}n + Temperatura + Voltaje $\times 6$ motores = \textbf{19 features} &
  El AAKR reconstruye la se\~{n}al cruda; necesita el tiempo y la posici\'{o}n para reconstruir la trayectoria \\
\bottomrule
\end{tabularx}
\end{center}

\subsection{Ventana Deslizante para Clasificadores}

Los clasificadores no trabajan punto a punto sino sobre \textit{ventanas temporales}:

\begin{lstlisting}[caption={Configuraci\'{o}n de ventana deslizante}]
window_size  = 50   # Numero de pasos temporales por ventana (5 segundos)
sample_step  = 20   # Desplazamiento entre ventanas consecutivas (2 segundos)
\end{lstlisting}

\begin{porquebox}
\textbf{Por qu\'{e} ventana deslizante en lugar de clasificaci\'{o}n punto a punto}\\[4pt]
\textbf{Problema de la clasificaci\'{o}n puntual}: una \'{u}nica muestra de sensor puede
ser ruidosa. Una lectura anormal en un instante no necesariamente indica un fallo;
puede ser un artefacto de transmisi\'{o}n o una fluctuaci\'{o}n normal. Decisiones basadas
en un solo punto tienen alta tasa de falsas alarmas.\\[4pt]
\textbf{Ventaja de la ventana de 50 muestras (5 segundos)}:
\begin{itemize}[itemsep=2pt]
  \item Agrega estad\'{i}sticas (media, std, tendencia) sobre un intervalo de 5 segundos.
  \item El modelo ve si la temperatura lleva 5 segundos subiendo (patr\'{o}n de fallo)
    o simplemente tuvo un pico puntual (ruido).
  \item Las estad\'{i}sticas de ventana son m\'{a}s robustas al ruido que los valores
    puntuales.
\end{itemize}
\textbf{Por qu\'{e} sample\_step = 20 (saltar 2 segundos entre ventanas)}:
\begin{itemize}[itemsep=2pt]
  \item Ventanas consecutivas que se solapan en el 90\% de sus muestras son
    pr\'{a}cticamente id\'{e}nticas: a\~{n}adir esas muestras solo a\~{n}ade redundancia.
  \item Step=20 reduce el n\'{u}mero de muestras en un factor 20, acelerando
    el entrenamiento y la CV sin perder informaci\'{o}n significativa.
  \item El solapamiento restante (30 muestras de 50) a\'{u}n garantiza que ning\'{u}n
    evento breve (de 2--3 s) cae completamente entre dos ventanas.
\end{itemize}
\end{porquebox}

% ============================================================
\section{Protocolo de Validaci\'{o}n Cruzada: Dise\~{n}o y Justificaci\'{o}n}
% ============================================================

\subsection{Por qu\'{e} Validaci\'{o}n Cruzada y no Partici\'{o}n Simple Train/Val}

En un dataset con 23 secuencias, una partici\'{o}n fija 80/20 dejar\'{i}a solo 4--5
secuencias para validaci\'{o}n. Si esas secuencias no contienen fallos (algo probable
dado que pocos experimentos tienen fallos), la estimaci\'{o}n del rendimiento ser\'{i}a
completamente sesgada. La validaci\'{o}n cruzada con $k=5$ garantiza que \textit{todas}
las secuencias aparecen exactamente una vez en validaci\'{o}n.

\begin{porquebox}
\textbf{Por qu\'{e} 5-fold cross-validation con k=5}\\[4pt]
El trade-off bias-varianza en la elecci\'{o}n de $k$:
\begin{itemize}[itemsep=2pt]
  \item \textbf{k peque\~{n}o (k=2)}: cada pliegue de train tiene el 50\% de los datos.
    El modelo est\'{a} subentrenado respecto al modelo final (que usar\'{a} el 100\%). La
    estimaci\'{o}n tiene \textbf{alto sesgo} (el modelo de validaci\'{o}n es peor que el real).
  \item \textbf{k grande (k=23, leave-one-out)}: cada pliegue de train tiene el 95\%
    de los datos. La estimaci\'{o}n tiene bajo sesgo pero \textbf{alta varianza}
    (los pliegues de validaci\'{o}n son muy peque\~{n}os: 1 secuencia cada vez).
  \item \textbf{k=5 (elegido)}: compromiso cl\'{a}sico. Cada fold de train tiene el 80\%
    de los datos ($\approx 18$ secuencias). La estimaci\'{o}n tiene varianza manejable.
\end{itemize}
\end{porquebox}

\subsection{Partici\'{o}n por Secuencia, no por Muestra}

\begin{advertenciabox}
\textbf{Por qu\'{e} NO particionar por muestra individual (data leakage temporal)}\\[4pt]
Si las 39\,309 filas se dividen aleatoriamente en folds, muestras temporalmente
cercanas de la misma secuencia aparecer\'{a}n en train Y en val simult\'{a}neamente.
Por ejemplo, la muestra del instante $t=100$ estar\'{a} en train y la del instante
$t=101$ (de la misma secuencia) estar\'{a} en val. El modelo aprende impl\'{i}citamente
la secuencia de test.\\[4pt]
Esto produce una sobreestimaci\'{o}n masiva del rendimiento. En series temporales
industriales, la partici\'{o}n correcta es siempre \textbf{a nivel de secuencia
experimental completa}, no a nivel de muestra.
\end{advertenciabox}

\subsection{B\'{u}squeda de Hiperpar\'{a}metros Anidada (Nested CV)}

Para RF y GB se utiliza \textit{nested cross-validation}: dentro de cada pliegue
externo ($k_{\text{ext}}=5$), se ejecuta un \texttt{GridSearchCV} interno de 3
pliegues sobre el conjunto de entrenamiento del pliegue externo.

\begin{center}
\begin{verbatim}
Fold externo k=1,2,3,4,5
├── Train externo (80%, ~18 secuencias)
│   ├── Fold interno k=1,2,3
│   │   ├── Train interno (67%, ~12 secuencias): fit del modelo
│   │   └── Val interno (33%, ~6 secuencias): score del modelo
│   └── GridSearch selecciona mejores hiperparámetros
├── Test externo (20%, ~5 secuencias): evaluacion FINAL del fold
│   └── Score con los mejores hiperparámetros seleccionados
└── Este score NO vio los datos de test externo durante la busqueda
\end{verbatim}
\end{center}

\begin{porquebox}
\textbf{Por qu\'{e} nested CV en lugar de CV simple con GridSearch}\\[4pt]
\textbf{Sin nested CV (CV plano)}: se realiza GridSearch sobre todos los datos de
train, se seleccionan los mejores hiperpar\'{a}metros, y luego se eval\'{u}a en el
fold de validaci\'{o}n. Problema: los hiperpar\'{a}metros se optimizaron teniendo en
cuenta (indirectamente) el fold de validaci\'{o}n, introduciendo \textit{optimism bias}.
La estimaci\'{o}n del rendimiento est\'{a} sesgada al alza.\\[4pt]
\textbf{Con nested CV}: los hiperpar\'{a}metros se seleccionan \'{u}nicamente en el
subconjunto de train del fold externo. El test externo es completamente desconocido
durante la b\'{u}squeda. La estimaci\'{o}n es no sesgada.\\[4pt]
\textbf{Costo computacional}: nested CV con 5 folds externos y 3 folds internos
requiere hasta $5 \times 3 \times N_{combos}$ entrenamientos. Para RF (12
combinaciones): $5 \times 3 \times 12 = 180$ fits. Para GB (24 combinaciones):
$5 \times 3 \times 24 = 360$ fits. Esto justifica el uso de \texttt{n\_jobs=-1}
para paralelizaci\'{o}n.
\end{porquebox}

\subsection{M\'{e}tricas de Evaluaci\'{o}n}

\begin{center}
\begin{tabular}{lp{5.5cm}p{4.5cm}}
\toprule
\textbf{M\'{e}trica} & \textbf{F\'{o}rmula} & \textbf{Justificaci\'{o}n} \\
\midrule
Accuracy & $\frac{TP+TN}{TP+TN+FP+FN}$ & Engañosa con clases desbalanceadas: un modelo que predice siempre ``normal'' alcanza 95\%+ de accuracy \\
\addlinespace
Precision & $\frac{TP}{TP+FP}$ & Fracción de alarmas que son fallos reales: crítica para evitar paradas innecesarias de producción \\
\addlinespace
Recall & $\frac{TP}{TP+FN}$ & Fracción de fallos detectados: crítica para seguridad; un fallo no detectado puede causar accidente \\
\addlinespace
\textbf{F1-score} & $\frac{2 \cdot P \cdot R}{P+R}$ & Media arm\'{o}nica de P y R; penaliza tanto el exceso de FP como el exceso de FN; \textbf{m\'{e}trica principal} \\
\bottomrule
\end{tabular}
\end{center}

\begin{porquebox}
\textbf{Por qu\'{e} F1 y no AUC-ROC como m\'{e}trica principal}\\[4pt]
El AUC-ROC es excelente para evaluar el poder discriminativo de un modelo en
\textit{cualquier umbral}. Sin embargo, en este problema:
\begin{itemize}[itemsep=2pt]
  \item La predicci\'{o}n final en Kaggle requiere una \textbf{decisi\'{o}n binaria} (0 o 1),
    no una probabilidad. El F1 eval\'{u}a directamente esa decisi\'{o}n.
  \item La m\'{e}trica p\'{u}blica de Kaggle es el \textbf{F1 promediado} sobre los 6 motores.
  \item El AUC-ROC puede ser muy optimista con clases extremadamente desbalanceadas
    (200:1) porque la curva ROC se basa en la tasa de verdaderos positivos y
    la tasa de falsos positivos, no en sus valores absolutos.
\end{itemize}
\end{porquebox}

% ============================================================
\section{Modelos Evaluados: Descripci\'{o}n y Justificaci\'{o}n de la Elecci\'{o}n}
% ============================================================

\subsection{Por qu\'{e} Cuatro Modelos Tan Diferentes}

\begin{porquebox}
\textbf{Racional de la selecci\'{o}n del portfolio de modelos}\\[4pt]
Se han elegido intencionalmente modelos que cubren el espacio de paradigmas de
aprendizaje relevantes para este problema:
\begin{itemize}[itemsep=3pt]
  \item \textbf{AAKR}: paradigma no supervisado/semisupervisado. Solo requiere datos
    normales para entrenar. Relevante porque en sistemas industriales nuevos, los
    datos de fallo son escasos o inexistentes al inicio.
  \item \textbf{RF}: paradigma supervisado no lineal, robusto. Baseline s\'{o}lido para
    datos tabulares con desbalanceo. Sirve para verificar que el problema es
    abordable con enfoques est\'{a}ndar.
  \item \textbf{GB (histogramado)}: paradigma supervisado no lineal, potencialmente
    m\'{a}s preciso que RF en tabular con tuning correcto. Sensible a hiperpar\'{a}metros.
  \item \textbf{LR}: paradigma supervisado lineal. Sirve como \textbf{l\'{i}nea base
    m\'{i}nima}: si el problema no es linealmente separable, LR falla y los otros
    deber\'{i}an superarla claramente.
\end{itemize}
El portfolio permite distinguir si el problema es lineal vs no lineal, y si la
informaci\'{o}n de fallo es suficiente para el aprendizaje supervisado.
\end{porquebox}

% ---- AAKR ----
\subsection{Modelo 1: AAKR (Auto-Associative Kernel Regression)}

\subsubsection{Descripci\'{o}n del Algoritmo}

El AAKR es un enfoque de \textit{novelty detection} basado en regresi\'{o}n. La idea
central es entrenar un regresor \textbf{exclusivamente con datos normales}: el modelo
aprende el comportamiento esperado del sensor bajo condiciones normales de operaci\'{o}n.
Durante la inferencia, se compara la temperatura predicha con la temperatura real;
una discrepancia superior a un umbral indica un fallo.

\begin{equation}
  \hat{T}^{(m)}_t = f_{\theta}(\mathbf{x}_t), \quad
  r_t = |T^{(m)}_t - \hat{T}^{(m)}_t|
\end{equation}

donde $r_t$ es el residuo en el instante $t$ y $m$ es el \'{i}ndice del motor.

La l\'{o}gica de detecci\'{o}n de fallos basada en residuos:

\begin{center}
\begin{verbatim}
Para cada instante t:
    Calcular residuo r_t = |T_real - T_predicha|
    Si r_t > threshold:
        marcar t como "anomalo"
    Si hay >=abnormal_limit instantes "anomalos"
    dentro de una ventana de window_size pasos:
        declarar FALLO en esa ventana
\end{verbatim}
\end{center}

\subsubsection{Implementaci\'{o}n}

\begin{lstlisting}[caption={Configuraci\'{o}n del detector AAKR}]
from sklearn.pipeline import Pipeline
from sklearn.preprocessing import StandardScaler
from sklearn.linear_model import LinearRegression

# Pipeline: normalizacion + regresion lineal
pipeline = Pipeline([
    ('scaler', StandardScaler()),
    ('regressor', LinearRegression())
])

# Wrapper de deteccion de fallos
detector = FaultDetectReg(
    model          = pipeline,
    threshold      = 0.9,   # Umbral de residuo normalizado
    abnormal_limit = 3,     # Minimo pasos consecutivos anomalos
    window_size    = 10,    # Ventana de suavizado del residuo
    sample_step    = 1      # Paso de muestreo (punto a punto)
)
\end{lstlisting}

\subsubsection{Ventajas, Inconvenientes y Contexto de Uso}

\begin{center}
\begin{tabular}{p{6cm}p{6cm}}
\toprule
\textbf{Ventajas} & \textbf{Inconvenientes} \\
\midrule
No requiere ejemplos de fallo para entrenar (aprendizaje semisupervisado) &
Sensible a la elecci\'{o}n del umbral $\tau$: un umbral bajo genera muchas
falsas alarmas; uno alto deja pasar fallos \\
\addlinespace
Interpretable: el residuo indica la magnitud de la anomal\'{i}a &
Un regresor lineal puede no capturar din\'{a}micas no lineales del sensor \\
\addlinespace
Generaliza bien a tipos de fallo no vistos en entrenamiento &
Requiere calibraci\'{o}n del umbral con datos de validaci\'{o}n \\
\addlinespace
\'{U}til al inicio del ciclo de vida del sistema, cuando los datos de fallo son escasos &
El recall alto pero la precisi\'{o}n baja observados indican demasiadas falsas alarmas \\
\bottomrule
\end{tabular}
\end{center}

% ---- RF ----
\subsection{Modelo 2: Random Forest (RF)}

\subsubsection{Descripci\'{o}n del Algoritmo}

Random Forest es un m\'{e}todo de \textit{ensemble} que combina m\'{u}ltiples \'{a}rboles de
decisi\'{o}n entrenados en submuestras bootstrap del conjunto de entrenamiento. La
predicci\'{o}n final se obtiene por votaci\'{o}n de mayor\'{i}a entre los $B = 200$ \'{a}rboles:

\begin{equation}
  \hat{y} = \text{argmax}_c \sum_{b=1}^{B} \mathbf{1}[h_b(\mathbf{x}) = c]
\end{equation}

El par\'{a}metro \texttt{class\_weight=`balanced'} ajusta autom\'{a}ticamente los pesos:

\begin{equation}
  w_c = \frac{N}{K \cdot N_c}
\end{equation}

donde $N$ es el n\'{u}mero total de muestras, $K$ el n\'{u}mero de clases y $N_c$ el
n\'{u}mero de muestras de la clase $c$.

\begin{porquebox}
\textbf{Por qu\'{e} class\_weight=`balanced' es crucial para este problema}\\[4pt]
Con una tasa de fallo del 0.5\% (Motor 3), hay aproximadamente 200 muestras normales
por cada muestra de fallo. Sin class\_weight, el \'{a}rbol aprende que ``predecir siempre
Normal'' tiene un error de entrenamiento del 0.5\%, lo cual es excelente desde el
punto de vista de la p\'{e}rdida de entropía cruzada. Nunca aprende a detectar fallos.\\[4pt]
Con \texttt{class\_weight=`balanced'}, cada muestra de fallo tiene un peso de
$w_1 = N/(K \cdot N_1) = 200$ veces el peso de una muestra normal.
Desde la perspectiva del optimizador, \textbf{los 200 ejemplos de fallo ``valen''
tanto como los 40\,000 ejemplos normales}. El modelo est\'{a} forzado a aprender
la frontera de decisi\'{o}n para fallos.
\end{porquebox}

\subsubsection{Configuraci\'{o}n e Hiperpar\'{a}metros}

\begin{lstlisting}[caption={Configuraci\'{o}n de Random Forest con GridSearchCV}]
from sklearn.ensemble import RandomForestClassifier
from sklearn.model_selection import GridSearchCV

# Modelo base
rf = RandomForestClassifier(
    class_weight  = 'balanced',
    n_estimators  = 200,
    random_state  = 42
)

# Grid de hiperparametros
param_grid_rf = {
    'max_depth'        : [5, 10, 20, None],
    'min_samples_leaf' : [1, 2, 5]
}
# 4 valores x 3 valores = 12 combinaciones
# 12 combinaciones x 3 folds internos = 36 fits por fold externo

gs_rf = GridSearchCV(
    rf, param_grid_rf,
    cv=3, scoring='f1', n_jobs=-1
)
\end{lstlisting}

% ---- GB ----
\subsection{Modelo 3: Gradient Boosting Histogramado (GB)}

\subsubsection{Descripci\'{o}n del Algoritmo}

\texttt{HistGradientBoostingClassifier} implementa gradient boosting sobre histogramas,
inspirado en LightGBM:

\begin{enumerate}[itemsep=2pt]
  \item \textbf{Paso 0}: Inicializar $F_0(\mathbf{x}) = \text{argmin}_c \mathcal{L}(y, c)$.
  \item \textbf{Paso k}: Calcular residuos pseudo-negativos: $r_i = -\partial \mathcal{L}/\partial F$.
  \item \textbf{Paso k}: Discretizar features en $B$ bins (histogramas) y ajustar \'{a}rbol $h_k$.
  \item \textbf{Actualizar}: $F_k = F_{k-1} + \eta \cdot h_k$, donde $\eta$ es el \texttt{learning\_rate}.
\end{enumerate}

\begin{porquebox}
\textbf{Por qu\'{e} GB puede colapsar con clases muy desbalanceadas}\\[4pt]
El Gradient Boosting con \texttt{learning\_rate} peque\~{n}o (0.01) y pocas muestras
de fallo es vulnerable a un fen\'{o}meno de ``congelaci\'{o}n'': los residuos de la clase
minoritaria son tan peque\~{n}os en magnitud (porque hay tan pocas) que el siguiente
\'{a}rbol en el boosting no aprende a corregirlos. El resultado es que el modelo
converge al equilibrio de ``predecir siempre 0'', especialmente cuando:
\begin{itemize}[itemsep=2pt]
  \item El \texttt{learning\_rate} es peque\~{n}o (0.01): cada \'{a}rbol contribuye poco.
  \item El n\'{u}mero de fallos es muy bajo ($<1\%$): la se\~{n}al de gradiente de la
    clase positiva es dominada por el gradiente de la clase negativa.
  \item El Motor 2 (17\% de fallos) deber\'{i}a ser m\'{a}s estable, pero a\'{u}n as\'{i}
    colapsa si el \texttt{learning\_rate=0.01} seleccionado por GridSearch en un fold
    particular resulta ser suboptimo para ese fold.
\end{itemize}
\end{porquebox}

\subsubsection{Configuraci\'{o}n e Hiperpar\'{a}metros}

\begin{lstlisting}[caption={Configuraci\'{o}n de Gradient Boosting con GridSearchCV}]
from sklearn.ensemble import HistGradientBoostingClassifier

# Modelo base
gb = HistGradientBoostingClassifier(
    class_weight = 'balanced',
    random_state = 42
)

# Grid de hiperparametros
param_grid_gb = {
    'max_iter'        : [100, 200],
    'learning_rate'   : [0.01, 0.1],
    'max_depth'       : [5, 10, None],
    'min_samples_leaf': [2, 5]
}
# 2 x 2 x 3 x 2 = 24 combinaciones
# 24 combinaciones x 3 folds internos = 72 fits por fold externo

gs_gb = GridSearchCV(
    gb, param_grid_gb,
    cv=3, scoring='f1', n_jobs=-1
)
\end{lstlisting}

% ---- LR ----
\subsection{Modelo 4: Regresi\'{o}n Log\'{i}stica (LR)}

\subsubsection{Descripci\'{o}n del Algoritmo}

La Regresi\'{o}n Log\'{i}stica con regularizaci\'{o}n \textit{elastic-net} combina L1 y L2:

\begin{equation}
  \min_{\mathbf{w}} \left[
    -\sum_{i} \left( y_i \log p_i + (1-y_i)\log(1-p_i) \right)
    + \frac{1}{C} \left( \rho \|\mathbf{w}\|_1 + \frac{1-\rho}{2}\|\mathbf{w}\|_2^2 \right)
  \right]
\end{equation}

donde $C$ es el inverso de la fuerza de regularizaci\'{o}n y $\rho$ es \texttt{l1\_ratio}.

\begin{porquebox}
\textbf{Por qu\'{e} LR falla en este problema}\\[4pt]
Las diferencias de se\~{n}al (\texttt{temp\_diff}, \texttt{volt\_diff}) que constituyen
las features principales no tienen una relaci\'{o}n lineal con el fallo en todos los
motores. Un fallo el\'{e}ctrico puede manifestarse como:
\begin{itemize}[itemsep=2pt]
  \item Voltaje bajo Y temperatura alta (interacci\'{o}n: no lineal).
  \item Voltaje bajo con temperatura normal (solo voltaje).
  \item Temperatura muy alta con voltaje normal (solo temperatura).
\end{itemize}
LR solo puede aprender ``voltaje bajo implica fallo'' o ``temperatura alta implica fallo''
por separado. No puede aprender la combinaci\'{o}n. RF y GB aprenden estas interacciones
en los nodos de los \'{a}rboles (``si voltaje$<$X Y temperatura$>$Y, entonces fallo'').
\end{porquebox}

\subsubsection{Configuraci\'{o}n e Hiperpar\'{a}metros}

\begin{lstlisting}[caption={Configuraci\'{o}n de Regresi\'{o}n Log\'{i}stica}]
from sklearn.linear_model import LogisticRegression

lr = LogisticRegression(
    penalty      = 'elasticnet',
    solver       = 'saga',
    class_weight = 'balanced',
    max_iter     = 300
)

param_grid_lr = {
    'C'        : [1, 10],
    'l1_ratio' : [0.5]
}
# 2 combinaciones x 3 folds = 6 fits por fold externo
\end{lstlisting}

% ============================================================
\section{Tabla de Hiperpar\'{a}metros: Espacio Buscado vs Mejor Valor}
% ============================================================

La siguiente tabla consolida el espacio de b\'{u}squeda de hiperpar\'{a}metros y los
mejores valores encontrados por GridSearchCV para los motores con resultados
disponibles:

\begin{center}
\begin{tabular}{llll}
\toprule
\textbf{Modelo} & \textbf{Hiperpar\'{a}metro} & \textbf{Rango buscado} & \textbf{Mejor valor (Motor 1)} \\
\midrule
\rowcolor{rowgray}
RF & \texttt{max\_depth} & $\{5,\, 10,\, 20,\, \text{None}\}$ & 10 \\
RF & \texttt{min\_samples\_leaf} & $\{1,\, 2,\, 5\}$ & 1 \\
\rowcolor{rowgray}
GB & \texttt{max\_iter} & $\{100,\, 200\}$ & 200 \\
GB & \texttt{learning\_rate} & $\{0.01,\, 0.1\}$ & 0.1 \\
\rowcolor{rowgray}
GB & \texttt{max\_depth} & $\{5,\, 10,\, \text{None}\}$ & 5 \\
GB & \texttt{min\_samples\_leaf} & $\{2,\, 5\}$ & 5 \\
\rowcolor{rowgray}
LR & \texttt{C} & $\{1,\, 10\}$ & 1 \\
LR & \texttt{l1\_ratio} & $\{0.5\}$ & 0.5 \\
\bottomrule
\end{tabular}
\end{center}

\begin{porquebox}
\textbf{Interpretaci\'{o}n de los mejores hiperpar\'{a}metros del Motor 1}\\[4pt]
\textbf{RF: max\_depth=10}. La profundidad 10 es intermedia: no tan profunda como para
sobreajustar (profundidad None genera \'{a}rboles de hasta 30--40 niveles con 30 features),
ni tan superficial como para perder complejidad. Con 30 features, profundidad 10
permite hasta $2^{10} = 1024$ regiones de decisi\'{o}n, suficiente para capturar
interacciones complejas.\\[4pt]
\textbf{GB: learning\_rate=0.1}. El GridSearch descart\'{o} learning\_rate=0.01, que era
probablemente demasiado lento para converger en 100--200 iteraciones con el desbalanceo
de este dataset.\\[4pt]
\textbf{GB: max\_depth=5}. \'{A}rboles m\'{a}s superficiales en GB que en RF es coherente
con la teor\'{i}a: el boosting necesita \'{a}rboles ``d\'{e}biles'' (shallow) para que el
ensamble global no sobreajuste al agregar muchos \'{a}rboles profundos.\\[4pt]
\textbf{LR: C=1}. La regularizaci\'{o}n fuerte ($C=1$, inverso de $\lambda$) indica que
el modelo tiende a sobreajustar con poca regularizaci\'{o}n ($C=10$). Con solo 23
secuencias de entrenamiento y 30 features, la regularizaci\'{o}n es necesaria.
\end{porquebox}

\subsection{N\'{u}mero Total de Entrenamientos por Motor}

\begin{center}
\begin{tabular}{lccc}
\toprule
\textbf{Modelo} & \textbf{Combos grid} & \textbf{Folds internos} & \textbf{Fits/fold externo} \\
\midrule
AAKR & -- (sin GridSearch) & -- & 1 \\
RF   & 12 & 3 & 36 \\
GB   & 24 & 3 & 72 \\
LR   & 2  & 3 & 6  \\
\midrule
\textbf{Total / fold externo} & & & \textbf{115} \\
\textbf{Total / motor (5 folds externos)} & & & \textbf{575} \\
\textbf{Total (6 motores)} & & & \textbf{3\,450} \\
\bottomrule
\end{tabular}
\end{center}

% ============================================================
\section{Resultados de la Validaci\'{o}n Cruzada}
% ============================================================

\subsection{Motor 1 (Tasa de fallo: $\approx 4$--$5\%$)}

\begin{center}
\begin{tabular}{lcccc}
\toprule
\textbf{Modelo} & \textbf{Accuracy} & \textbf{Precision} & \textbf{Recall} & \textbf{F1-score} \\
\midrule
AAKR & 0.7112 & 0.0719 & 0.3825 & 0.1204 \\
\textbf{RF}   & \textbf{0.9188} & \textbf{0.4000} & \textbf{0.4000} & \textbf{0.4000} \\
GB   & 0.8426 & 0.2181 & 0.2788 & 0.2295 \\
LR   & 0.7441 & 0.0540 & 0.1642 & 0.0705 \\
\bottomrule
\end{tabular}
\end{center}

\textbf{An\'{a}lisis}: RF obtiene el mejor F1 (0.40). Notable que Precision = Recall = 0.40
exactamente para RF: esto ocurre cuando el modelo tiene igual n\'{u}mero de FP y FN,
lo que indica que el umbral de decisi\'{o}n por defecto (0.5) es aproximadamente optimal
para este motor. AAKR tiene recall 0.38 pero precisi\'{o}n 0.07, indicando 5 falsas
alarmas por cada fallo real detectado.

\subsection{Motor 2 (Tasa de fallo: $\approx 17\%$)}

\begin{center}
\begin{tabular}{lcccc}
\toprule
\textbf{Modelo} & \textbf{Accuracy} & \textbf{Precision} & \textbf{Recall} & \textbf{F1-score} \\
\midrule
AAKR & 0.6913 & 0.1759 & 0.2888 & 0.1540 \\
\textbf{RF}   & \textbf{0.8657} & \textbf{0.4000} & \textbf{0.4000} & \textbf{0.4000} \\
GB   & 0.8597 & 0.0269 & 0.0004 & 0.0008 \\
LR   & 0.7455 & 0.0011 & 0.0009 & 0.0010 \\
\bottomrule
\end{tabular}
\end{center}

\textbf{An\'{a}lisis}: El Motor 2 tiene la mayor tasa de fallos (17\%), lo que en principio
facilita el aprendizaje. Sin embargo, GB y LR colapsan a F1 $\approx 0$, confirmando
que las features no son linealmente separables para este motor y que GB es inestable
en este r\'{e}gimen. La accuracy de GB (0.86) es casi id\'{e}ntica a RF (0.87), lo que
significa que GB tambi\'{e}n predice casi siempre ``Normal'' y ``acierta'' gracias a la
clase mayoritaria. El verdadero indicador de fallo es el F1: GB = 0.0008 vs RF = 0.40.

\subsection{Motor 3 (Tasa de fallo: $< 0.5\%$)}

\begin{center}
\begin{tabular}{lcccc}
\toprule
\textbf{Modelo} & \textbf{Accuracy} & \textbf{Precision} & \textbf{Recall} & \textbf{F1-score} \\
\midrule
AAKR & 0.8488 & 0.0215 & 0.1663 & 0.0381 \\
\textbf{RF}   & \textbf{0.9949} & \textbf{0.6000} & \textbf{0.6000} & \textbf{0.6000} \\
GB   & 0.9592 & 0.0000 & 0.0000 & 0.0000 \\
LR   & 0.6846 & 0.0000 & 0.0000 & 0.0000 \\
\bottomrule
\end{tabular}
\end{center}

\textbf{An\'{a}lisis}: El Motor 3 es el caso m\'{a}s extremo (200:1). GB y LR obtienen F1 = 0
(el modelo predice siempre ``Normal''). Paradoja: GB tiene accuracy 0.9592 pero F1=0,
lo que ilustra perfectamente por qu\'{e} la accuracy es una m\'{e}trica engañosa en problemas
desbalanceados. RF con F1=0.60 es un resultado notable, logrado \'{u}nicamente gracias
a \texttt{class\_weight=`balanced'}.

\subsection{Motores 4, 5 y 6 (Error de Ejecuci\'{o}n)}

\begin{center}
\begin{tabular}{lccccc}
\toprule
\textbf{Motor} & \textbf{AAKR} & \textbf{RF} & \textbf{GB} & \textbf{LR} & \textbf{Causa} \\
\midrule
Motor 4 & 0.55 (acc.) & Error & Error & No ejecutado & NameError en kernel \\
Motor 5 & No ejecutado & Error & Error & No ejecutado & NameError en kernel \\
Motor 6 & No ejecutado & Error & Error & No ejecutado & NameError en kernel \\
\bottomrule
\end{tabular}
\end{center}

El error \texttt{NameError} ocurri\'{o} porque el notebook de Kaggle perdi\'{o} el estado
del kernel al reiniciar entre celdas, y algunas variables (probablemente los dataframes
preprocesados o las funciones de utilidad) no estaban definidas en el momento de la
ejecuci\'{o}n de las celdas de los Motores 4--6.

\subsection{Resumen Comparativo de los Tres Motores Evaluados}

\begin{center}
\begin{tabular}{lccccc}
\toprule
\textbf{Motor} & \textbf{Fallos (\%)} & \textbf{Mejor Modelo} & \textbf{F1} & \textbf{Accuracy} & \textbf{Observaci\'{o}n} \\
\midrule
Motor 1 & $\sim 4{-}5\%$  & RF & 0.40 & 0.92 & Ratio $\sim 20:1$, CV estable \\
Motor 2 & $\sim 17\%$     & RF & 0.40 & 0.87 & GB colapsa pese a m\'{a}s fallos \\
Motor 3 & $< 0.5\%$       & RF & 0.60 & 0.99 & Ratio $>200:1$; GB y LR a F1=0 \\
Motor 4 & $\sim 17\%$     & --  & --   & --   & Error de ejecuci\'{o}n (NameError) \\
Motor 5 & Desconocida     & --  & --   & --   & Error de ejecuci\'{o}n (NameError) \\
Motor 6 & Desconocida     & --  & --   & --   & Error de ejecuci\'{o}n (NameError) \\
\bottomrule
\end{tabular}
\end{center}

% ============================================================
\section{An\'{a}lisis Profundo de Resultados}
% ============================================================

\subsection{Por qu\'{e} RF es el Mejor Modelo: An\'{a}lisis Multifactorial}

\begin{porquebox}
\textbf{Explicaci\'{o}n completa del dominio de RF en este problema}\\[4pt]
\textbf{Factor 1: No-linealidad robusta}. Los l\'{i}mites de decisi\'{o}n de RF son
hiperplanos por piezas, capturando interacciones complejas entre sensores. La
detecci\'{o}n de fallos en motores industriales t\'{i}picamente requiere reglas
del tipo ``si voltaje < X AND temperatura > Y'' que LR no puede aprender.\\[4pt]
\textbf{Factor 2: Estabilidad del ensemble}. Con 200 \'{a}rboles, cada uno entrenado
en un bootstrap de $\approx 63\%$ de las muestras, la varianza del estimador es
$B$ veces menor que la de un solo \'{a}rbol ($B=200$). Con tan pocas muestras de
fallo, esta reducci\'{o}n de varianza es cr\'{i}tica para la estabilidad del F1.\\[4pt]
\textbf{Factor 3: class\_weight=`balanced' bien calibrado para RF}. El boosting (GB)
tambi\'{e}n tiene \texttt{class\_weight=`balanced'}, pero el mecanismo es diferente:
en GB, el peso se aplica a los residuos en cada iteraci\'{o}n de boosting. Si la tasa
de aprendizaje es peque\~{n}a, los residuos ponderados de la clase minoritaria convergen
lentamente y el modelo puede quedar atrapado en un m\'{i}nimo local donde predice siempre
0. En RF, el peso se aplica directamente al criterio de divisi\'{o}n de cada nodo,
lo que es m\'{a}s estable.\\[4pt]
\textbf{Factor 4: Regularizaci\'{o}n impl\'{i}cita por bootstrap y feature bagging}.
Incluso sin max\_depth limitado, los \'{a}rboles individuales de RF est\'{a}n regularizados
por el hecho de ver solo $\sqrt{30} \approx 5$ features aleatorias por nodo. Esto
reduce la correlaci\'{o}n entre \'{a}rboles y mejora la generalizaci\'{o}n.
\end{porquebox}

\subsection{Alta Varianza entre Pliegues: Un Fen\'{o}meno Esperado}

Un resultado sorprendente a primera vista es la alta varianza del F1 entre pliegues:
en algunos pliegues F1 = 1.0 y en otros F1 = 0.0. Esta secci\'{o}n explica por qu\'{e}
esto es un resultado \textbf{esperado y correcto}, no un error de implementaci\'{o}n.

\begin{porquebox}
\textbf{Modelo estad\'{i}stico de la varianza entre folds}\\[4pt]
Sea $n_f = 23$ el n\'{u}mero de secuencias. Supongamos que los fallos se concentran
en $n_+ = 2$ secuencias (el caso extremo observado). Con $k=5$ folds:
\begin{itemize}[itemsep=2pt]
  \item Cada fold de validaci\'{o}n contiene $23/5 \approx 4.6$ secuencias.
  \item La probabilidad de que \textbf{ninguna} secuencia con fallos cai\-ga en un fold
    de validaci\'{o}n es: $P(\text{ninguna fallo en val}) \approx \binom{21}{4.6}/\binom{23}{4.6}$.
  \item Cuando no hay fallos en validaci\'{o}n: F1 = 0.0 por definici\'{o}n (no hay positivos
    verdaderos; el denominador del recall es 0).
  \item Cuando hay una secuencia con fallos en validaci\'{o}n y el modelo la detecta
    correctamente: F1 puede llegar a 1.0.
\end{itemize}
\textbf{Conclusi\'{o}n}: La varianza entre folds NO indica inestabilidad del modelo;
indica que el dataset tiene pocas secuencias con fallos. El promedio del F1 entre
folds sigue siendo una estimaci\'{o}n no sesgada del rendimiento esperado, pero su
desviaci\'{o}n est\'{a}ndar es inevitablemente alta.
\end{porquebox}

\subsection{Tabla de An\'{a}lisis de Varianza entre Folds}

La siguiente tabla ilustra c\'{o}mo la distribuci\'{o}n de fallos en los folds afecta
el F1 observado. Los valores son ilustrativos basados en el comportamiento observado:

\begin{center}
\begin{tabular}{ccccc}
\toprule
\textbf{Fold} & \textbf{Secuencias val} & \textbf{Con fallos en val} & \textbf{F1 RF (aprox.)} & \textbf{F1 AAKR (aprox.)} \\
\midrule
Fold 1 & $\sim 5$ & 0 & 0.00 & 0.00 \\
Fold 2 & $\sim 5$ & 1 & 0.40--0.80 & 0.10--0.20 \\
Fold 3 & $\sim 5$ & 0 & 0.00 & 0.00 \\
Fold 4 & $\sim 4$ & 2 & 0.60--1.00 & 0.15--0.30 \\
Fold 5 & $\sim 4$ & 0 & 0.00 & 0.00 \\
\midrule
\textbf{Media} & & & \textbf{0.20--0.40} & \textbf{0.05--0.10} \\
\textbf{Std} & & & \textbf{alta} & \textbf{alta} \\
\bottomrule
\end{tabular}
\end{center}

\begin{advertenciabox}
\textbf{La desviaci\'{o}n est\'{a}ndar alta es \textbf{esperada} con pocas secuencias de fallo}\\[4pt]
Esta observaci\'{o}n tiene implicaciones importantes para el dise\~{n}o del sistema de
detecci\'{o}n en producci\'{o}n:
\begin{enumerate}[itemsep=2pt]
  \item El rendimiento \textbf{en producci\'{o}n} depender\'{a} de si los fallos futuros
    se parecen a los vistos en entrenamiento.
  \item Si el tipo de fallo es nuevo (no visto en train), el modelo puede tener
    F1 cercano a 0.0, independientemente de su rendimiento en CV.
  \item Esto justifica implementar AAKR como sistema complementario: detecta
    anomal\'{i}as nuevas aunque no las haya visto en entrenamiento.
\end{enumerate}
\end{advertenciabox}

\subsection{Colapso de GB y LR: An\'{a}lisis Detallado}

\begin{center}
\begin{tabular}{p{3cm}p{4cm}p{4.5cm}}
\toprule
\textbf{Modelo} & \textbf{Mecanismo de colapso} & \textbf{Soluci\'{o}n potencial} \\
\midrule
GB (Motor 2, F1=0.0008) &
  Con \texttt{learning\_rate=0.1} y muchas iteraciones (200), el modelo puede
  sobreajustar a la distribuci\'{o}n del fold de train y generalizar mal a val &
  Reducir el espacio de b\'{u}squeda de learning\_rate; añadir \texttt{l2\_regularization}
  al grid; usar \texttt{early\_stopping} \\
\addlinespace
GB (Motor 3, F1=0) &
  Con 0.5\% de fallos, incluso con class\_weight=`balanced', la señal de gradiente
  de la clase positiva es demasiado peque\~{n}a relativa al ruido de optimizaci\'{o}n &
  Combinar SMOTE + GB; usar Isolation Forest como alternativa semisupervisada \\
\addlinespace
LR (todos los motores) &
  Las features (temp\_diff, volt\_diff) tienen interacciones no lineales con el
  estado de fallo que LR no puede capturar con un hiperplano &
  Añadir features polin\'{o}micas o de interacci\'{o}n; usar SVM con kernel RBF \\
\bottomrule
\end{tabular}
\end{center}

\subsection{Comparaci\'{o}n AAKR vs Clasificadores Supervisados}

El AAKR muestra un patr\'{o}n consistente: recall relativamente alto pero precisi\'{o}n
muy baja. La raz\'{o}n es la siguiente:

\begin{center}
\begin{verbatim}
AAKR: residuo r_t > threshold = "anomalia"
    - Threshold bajo (0.9 normalizado) -> muchas anomalias detectadas
    - Pero muchos estados "normales" tambien tienen residuos > 0.9
      (la distribucion de residuos normales tiene cola larga)
    - Resultado: recall alto, precision baja, muchas falsas alarmas

Clasificadores supervisados:
    - Aprenden directamente la frontera normal vs fallo
    - No necesitan calibrar un umbral sobre residuos
    - Pueden usar combinaciones de features para decidir
    - Resultado: mejor equilibrio precision/recall cuando hay datos de fallo
\end{verbatim}
\end{center}

\begin{claveBox}
\textbf{Lección clave: AAKR complementa, no reemplaza}\\[4pt]
AAKR es \'{u}til en la fase inicial de un sistema industrial (pocos datos de fallo)
y para detectar \textbf{tipos de fallo nuevos} (no vistos en entrenamiento). Los
clasificadores supervisados son superiores una vez que se tienen suficientes datos
etiquetados. En producci\'{o}n, un sistema robusto usar\'{i}a \textbf{ambos en paralelo}:
AAKR para nuevos tipos de fallo, RF para fallos conocidos.
\end{claveBox}

% ============================================================
\section{Submission a Kaggle: El Pipeline Completo}
% ============================================================

\subsection{Por qu\'{e} la Submission Eval\'{u}a el Pipeline Completo}

\begin{porquebox}
\textbf{Por qu\'{e} la submission a Kaggle es el verdadero test del pipeline}\\[4pt]
El conjunto de test de Kaggle (14\,157 filas, 8 secuencias) es completamente
desconocido durante el desarrollo. A diferencia de la validaci\'{o}n cruzada, donde
los pliegues de validaci\'{o}n son subconjuntos del training, el test de Kaggle puede
contener:
\begin{itemize}[itemsep=2pt]
  \item Condiciones operativas nuevas no vistas en entrenamiento.
  \item Tipos de fallo diferentes a los del entrenamiento.
  \item Diferentes distribuciones de temperatura ambiente o carga mec\'{a}nica.
\end{itemize}
La m\'{e}trica p\'{u}blica de Kaggle es el \textbf{F1 macro promediado sobre los 6 motores},
por lo que un motor con F1=0 penaliza fuertemente el resultado global, aunque los
otros cinco funcionen perfectamente.
\end{porquebox}

\subsection{Proceso de Predicci\'{o}n}

\begin{lstlisting}[caption={Proceso de entrenamiento final y predicci\'{o}n}]
import pandas as pd
from sklearn.ensemble import RandomForestClassifier

# 1. Cargar y preprocesar datos de test
df_test = pd.read_csv('test_data.csv')
pre_processing(df_test)   # Mismos pasos que en entrenamiento

# 2. Para cada motor, entrenar en TODOS los datos de entrenamiento
best_params = {
    1: {'max_depth': 10, 'min_samples_leaf': 1},
    2: {'max_depth': 10, 'min_samples_leaf': 1},
    3: {'max_depth': 10, 'min_samples_leaf': 1},
    # Motores 4, 5, 6: usar parametros de Motor 1 como fallback
    4: {'max_depth': 10, 'min_samples_leaf': 1},
    5: {'max_depth': 10, 'min_samples_leaf': 1},
    6: {'max_depth': 10, 'min_samples_leaf': 1},
}

predictions = {}
for motor_id in range(1, 7):
    X_train = extract_features(df_train, motor_id)
    y_train = df_train[f'data_motor_{motor_id}_label']

    best_rf = RandomForestClassifier(
        class_weight     = 'balanced',
        n_estimators     = 200,
        max_depth        = best_params[motor_id]['max_depth'],
        min_samples_leaf = best_params[motor_id]['min_samples_leaf'],
        random_state     = 42
    )
    best_rf.fit(X_train, y_train)

    X_test = extract_features(df_test, motor_id)
    predictions[f'data_motor_{motor_id}_label'] = best_rf.predict(X_test)

# 3. Construir dataframe de submission
submission = pd.DataFrame({
    'idx'               : df_test.index,
    'data_motor_1_label': predictions['data_motor_1_label'],
    'data_motor_2_label': predictions['data_motor_2_label'],
    'data_motor_3_label': predictions['data_motor_3_label'],
    'data_motor_4_label': predictions['data_motor_4_label'],
    'data_motor_5_label': predictions['data_motor_5_label'],
    'data_motor_6_label': predictions['data_motor_6_label'],
    'test_condition'    : df_test['test_condition']
})

submission.to_csv('submission.csv', index=False)
\end{lstlisting}

\begin{advertenciabox}
\textbf{Por qu\'{e} entrenar en el 100\% de los datos de entrenamiento para la submission}\\[4pt]
Durante la validaci\'{o}n cruzada se usaba el 80\% de los datos para entrenar el modelo
evaluado. Para la submission final, se usan los 39\,309 ejemplos completos. Esto es
correcto porque:
\begin{itemize}[itemsep=2pt]
  \item Ya no necesitamos estimar el rendimiento; lo estimamos con CV.
  \item Mais datos de entrenamiento $\Rightarrow$ mejor capacidad de generalizaci\'{o}n
    al conjunto de test de Kaggle.
  \item Con solo 23 secuencias, usar el 100\% en lugar del 80\% añade $\approx 5$
    secuencias adicionales, incluyendo potencialmente secuencias con fallos que
    enriquecen la representaci\'{o}n de la clase positiva.
\end{itemize}
\end{advertenciabox}

\subsection{Estructura del Archivo de Submission}

\begin{center}
\begin{tabular}{llp{5.5cm}}
\toprule
\textbf{Columna} & \textbf{Tipo} & \textbf{Descripci\'{o}n} \\
\midrule
\texttt{idx}                     & int   & \'{I}ndice de la fila en el dataset de test \\
\texttt{data\_motor\_1\_label}   & int   & Predicci\'{o}n Motor 1 (0=normal, 1=fallo) \\
\texttt{data\_motor\_2\_label}   & int   & Predicci\'{o}n Motor 2 (0=normal, 1=fallo) \\
\texttt{data\_motor\_3\_label}   & int   & Predicci\'{o}n Motor 3 (0=normal, 1=fallo) \\
\texttt{data\_motor\_4\_label}   & int   & Predicci\'{o}n Motor 4 (0=normal, 1=fallo) \\
\texttt{data\_motor\_5\_label}   & int   & Predicci\'{o}n Motor 5 (0=normal, 1=fallo) \\
\texttt{data\_motor\_6\_label}   & int   & Predicci\'{o}n Motor 6 (0=normal, 1=fallo) \\
\texttt{test\_condition}         & str   & Condici\'{o}n de test de la secuencia \\
\bottomrule
\end{tabular}
\end{center}

\subsection{Riesgo de Degradaci\'{o}n en Kaggle vs CV}

\begin{center}
\begin{tabular}{p{5cm}p{4.5cm}p{3.5cm}}
\toprule
\textbf{Fuente de degradaci\'{o}n} & \textbf{Efecto probable} & \textbf{Mitigaci\'{o}n} \\
\midrule
Tipos de fallo nuevos en test &
  F1 bajo o nulo para los motores afectados &
  Complementar con AAKR \\
\addlinespace
Distribución de temperatura/voltaje diferente en test &
  Las features diff pueden tener magnitudes diferentes &
  La compensaci\'{o}n de sesgo reduce este riesgo \\
\addlinespace
Motores 4--6 sin hiperpar\'{a}metros optimizados &
  Se usan hiperpar\'{a}metros del Motor 1 como proxy; puede ser suboptimal &
  Corregir NameError y ejecutar CV completa \\
\addlinespace
Alta proporci\'{o}n de fallos en test vs train &
  El modelo puede estar subcalibrado para la clase positiva &
  Ajustar umbral de decisi\'{o}n con Precision-Recall curve \\
\bottomrule
\end{tabular}
\end{center}

% ============================================================
\section{Discusi\'{o}n: Limitaciones y Alternativas}
% ============================================================

\subsection{Limitaciones del Enfoque Actual}

\subsubsection{Limitaci\'{o}n 1: Ventana Deslizante como \'{U}nico Mecanismo Temporal}

La ventana deslizante con estad\'{i}sticas simples (media, std) sobre 50 muestras es
una aproximaci\'{o}n eficiente pero que pierde informaci\'{o}n sobre el \textbf{orden temporal}
dentro de la ventana. Una temperatura que sube de 30 a 45\textdegree{}C en 5 segundos
tiene la misma media que una que baja de 45 a 30\textdegree{}C, pero el primero es
mucho m\'{a}s alarmanete. Modelos secuenciales (LSTM, Transformer) preservar\'{i}an
esta informaci\'{o}n de orden.

\subsubsection{Limitaci\'{o}n 2: Una \'{U}nica Capa de Diferencias}

Solo se calculan diferencias de primer orden (lag=10). La segunda derivada temporal
($\Delta^2 x_t = \Delta x_t - \Delta x_{t-10}$) capturar\'{i}a la \textbf{aceleraci\'{o}n}
del cambio de se\~{n}al, potencialmente m\'{a}s informativa para detectar la \textit{inminencia}
de un fallo (cuando la tasa de cambio en s\'{i} misma empieza a aumentar).

\subsubsection{Limitaci\'{o}n 3: No se Explotan las Correlaciones entre Motores}

Los seis motores operan en el mismo sistema rob\'{o}tico y est\'{a}n f\'{i}sicamente acoplados.
Un fallo en el Motor 1 puede afectar las cargas sobre el Motor 2. El modelo actual
trata cada motor de forma independiente, perdiendo esta informaci\'{o}n de acoplamiento.

\subsection{Alternativas T\'{e}cnicas para Mejoras Futuras}

\begin{center}
\begin{tabular}{p{3.5cm}p{5cm}p{4cm}}
\toprule
\textbf{T\'{e}cnica} & \textbf{Ventaja} & \textbf{Costo} \\
\midrule
SMOTE sobre features de ventana &
  Genera ejemplos sint\'{e}ticos en el espacio de features para la clase minoritaria &
  Riesgo de sobreajuste si el espacio sintético no representa bien la distribución real \\
\addlinespace
Features FFT (espectro de frecuencias) &
  Detecta componentes oscilatorias anómalas (vibraciones, frecuencias de fallo) &
  Mayor coste computacional; requiere señales estacionarias \\
\addlinespace
LSTM bidireccional &
  Captura dependencias temporales de largo alcance en ambas direcciones &
  Requiere mucho más dato y cómputo; riesgo de overfitting con 23 secuencias \\
\addlinespace
Isolation Forest &
  Detecci\'{o}n de anomalías sin labels de fallo; complementario al RF supervisado &
  El recall puede ser bajo si los fallos no son anomalías en el espacio de features \\
\addlinespace
Calibraci\'{o}n del umbral (Platt scaling) &
  Ajusta el umbral de decisi\'{o}n 0.5 a un valor \'{o}ptimo por motor &
  Requiere un conjunto de calibraci\'{o}n separado \\
\bottomrule
\end{tabular}
\end{center}

% ============================================================
\section{Conclusiones}
% ============================================================

\subsection{Modelo Seleccionado y Justificaci\'{o}n}

\textbf{Random Forest con \texttt{class\_weight=`balanced'}} es el mejor modelo para
la detecci\'{o}n de fallos en los 6 motores servo-rob\'{o}ticos, por los siguientes motivos:

\begin{enumerate}[itemsep=4pt]
  \item \textbf{Consistencia}: Es el \'{u}nico modelo que obtiene F1 $> 0$ en todos los
    motores evaluados, incluido el Motor 3 con $< 0.5\%$ de fallos (ratio 200:1).

  \item \textbf{Mejor F1-score}: F1 = 0.40 (Motores 1 y 2), F1 = 0.60 (Motor 3),
    frente a F1 $< 0.25$ del resto de modelos en los motores comparables.

  \item \textbf{Equilibrio precisi\'{o}n/recall}: El balance de clases integrado en
    \texttt{class\_weight=`balanced'} evita el sesgo hacia la clase mayoritaria
    de forma m\'{a}s estable que GB.

  \item \textbf{Robustez}: El promediado de 200 \'{a}rboles reduce la sensibilidad a
    secuencias de entrenamiento con pocos fallos, produciendo estimadores estables.

  \item \textbf{No-linealidad sin sobreajuste}: La combinaci\'{o}n de bootstrap,
    feature bagging y max\_depth=10 proporciona capacidad no lineal con
    regularizaci\'{o}n impl\'{i}cita.
\end{enumerate}

\subsection{Lecciones Aprendidas}

\begin{enumerate}[itemsep=4pt]
  \item \textbf{El balanceo de clases es la decisi\'{o}n m\'{a}s cr\'{i}tica}: Ante una
    asimetr\'{i}a de hasta 200:1, ning\'{u}n modelo --por sofisticado que sea-- puede
    detectar fallos sin mecanismos expl\'{i}citos de compensaci\'{o}n de clases.
    \texttt{class\_weight=`balanced'} es la herramienta m\'{a}s simple y efectiva.

  \item \textbf{La validaci\'{o}n cruzada por secuencia es obligatoria}: La partici\'{o}n
    por muestras individuales introduce data leakage temporal y sobreestima
    masivamente el rendimiento.

  \item \textbf{La alta varianza entre pliegues no es un defecto}: Es informaci\'{o}n
    valiosa sobre la distribuci\'{o}n real de fallos. Indica que los fallos est\'{a}n
    concentrados en pocas secuencias y que el sistema debe estar preparado para
    distribuciones de fallo cambiantes en producci\'{o}n.

  \item \textbf{El preprocessing importa tanto como el modelo}: Las transformaciones
    de compensaci\'{o}n de sesgo y c\'{a}lculo de diferencias (que convierten se\~{n}ales
    absolutas no estacionarias en desviaciones din\'{a}micas relativas) son
    fundamentales para que los modelos funcionen. Un RF sin estas transformaciones
    probablemente obtendr\'{i}a F1 cercano a 0.

  \item \textbf{AAKR y RF son complementarios}: AAKR detecta anomal\'{i}as nuevas;
    RF clasifica fallos conocidos. Un sistema de producci\'{o}n robusto combinar\'{i}a
    ambos enfoques.

  \item \textbf{El n\'{u}mero de secuencias limita la CV, no el n\'{u}mero de muestras}:
    Pasar de 39\,309 a 200\,000 muestras sin a\~{n}adir secuencias no mejora
    significativamente la estimaci\'{o}n del rendimiento por CV, ya que la varianza
    est\'{a} dominada por el n\'{u}mero de secuencias con fallos ($\approx 2$).
\end{enumerate}

\subsection{Trabajo Futuro}

\begin{itemize}[itemsep=3pt]
  \item \textbf{Corregir el NameError} para completar la evaluaci\'{o}n de los Motores
    4, 5 y 6 con todos los modelos supervisados.
  \item \textbf{SMOTE o datos sint\'{e}ticos}: Sobremuestrear la clase minoritaria en
    el espacio de features de ventana para proporcionar m\'{a}s ejemplos de fallo.
  \item \textbf{Features espectrales (FFT)}: Extraer componentes de frecuencia
    de las se\~{n}ales en cada ventana para detectar vibraciones peri\'{o}dicas de fallo.
  \item \textbf{Modelos secuenciales}: Explorar LSTM o Transformer para aprovechar la
    estructura temporal de orden de las se\~{n}ales dentro de la ventana.
  \item \textbf{Calibraci\'{o}n del umbral}: Optimizar el umbral de decisi\'{o}n por motor
    usando la curva precisi\'{o}n-recall, potencialmente mejorando el F1 en 0.05--0.15.
  \item \textbf{Ensemble de AAKR + RF}: Combinar las predicciones de AAKR y RF
    mediante un meta-clasificador para obtener lo mejor de ambos paradigmas.
\end{itemize}

% ============================================================
\appendix
\section*{Ap\'{e}ndice A: Resumen del Entorno de Desarrollo}
\addcontentsline{toc}{section}{Ap\'{e}ndice A: Entorno de Desarrollo}
% ============================================================

\begin{center}
\begin{tabular}{ll}
\toprule
\textbf{Componente} & \textbf{Versi\'{o}n/Valor} \\
\midrule
Python              & 3.10+ \\
scikit-learn        & 1.3+ \\
pandas              & 2.0+ \\
numpy               & 1.24+ \\
Plataforma          & Kaggle Notebook / Jupyter \\
Validaci\'{o}n cruzada  & \texttt{n\_cv = 5} (externo), \texttt{cv = 3} (GridSearch interno) \\
Ventana deslizante  & \texttt{window\_size = 50, sample\_step = 20} \\
Semilla aleatoria   & \texttt{random\_state = 42} \\
n\_estimators RF    & 200 \\
Frecuencia muestreo & 10 Hz (0.1 s/muestra) \\
Lag diferencias     & 10 pasos = 1 segundo \\
\bottomrule
\end{tabular}
\end{center}

% ============================================================
\section*{Ap\'{e}ndice B: Diagrama Conceptual del Pipeline Completo}
\addcontentsline{toc}{section}{Ap\'{e}ndice B: Pipeline Completo}
% ============================================================

\begin{center}
\begin{verbatim}
DATOS CRUDOS (39309 filas x 19 columnas)
        |
        | 1. remove_outliers(): clipping fisico + forward-fill
        v
DATOS LIMPIOS
        |
        | 2. compensate_seq_bias(): restar x[0] de cada secuencia
        v
DATOS CENTRADOS (señales relativas al estado inicial)
        |
        | 3. cal_diff(lag=10): añadir 18 columnas de diferencias
        v
DATOS CON FEATURES DINAMICAS (37 columnas totales)
        |
        +----------+----------+----------+
        |          |          |          |
        v          v          v          v
     AAKR         RF         GB         LR
  (19 feats)  (30 feats)  (30 feats) (30 feats)
     |           |           |          |
     | FaultDetectReg  VentanaDeslizante (w=50, s=20)
     v           v           v          v
  PREDICCION  PREDICCION  PREDICCION PREDICCION
  punto a     por ventana por ventana por ventana
  punto
        |          |          |          |
        +----------+----------+----------+
                   |
                   | GridSearchCV anidado (3 folds internos)
                   | para RF, GB y LR
                   v
            5-FOLD CV EXTERNO
            (por secuencia, no por muestra)
                   |
                   v
            F1, Precision, Recall, Accuracy
            por fold y por modelo
                   |
                   | Seleccion: mejor F1 medio
                   v
            MODELO GANADOR (RF en motores 1, 2, 3)
                   |
                   | Reentrenar en 100% training data
                   v
            PREDICCION EN TEST (14157 filas)
                   |
                   v
            submission.csv -> Kaggle
\end{verbatim}
\end{center}

% ============================================================
\section*{Ap\'{e}ndice C: Glosario de T\'{e}rminos}
\addcontentsline{toc}{section}{Ap\'{e}ndice C: Glosario}
% ============================================================

\begin{center}
\begin{tabularx}{\textwidth}{lX}
\toprule
\textbf{T\'{e}rmino} & \textbf{Definici\'{o}n} \\
\midrule
AAKR & Auto-Associative Kernel Regression: regresor entrenado solo con datos normales que detecta anomal\'{i}as por residuo de predicci\'{o}n. \\
\addlinespace
Nested CV & Validaci\'{o}n cruzada anidada: CV externo para evaluar, CV interno para seleccionar hiperpar\'{a}metros. Garantiza estimaci\'{o}n no sesgada. \\
\addlinespace
class\_weight=`balanced' & Ajuste autom\'{a}tico de pesos de clase: $w_c = N/(K \cdot N_c)$. Hace que la clase minoritaria tenga el mismo peso total que la mayoritaria. \\
\addlinespace
Data leakage & Fuga de informaci\'{o}n: cuando datos del conjunto de validaci\'{o}n/test influyen impl\'{i}citamente en el entrenamiento, produciendo sobreestimaci\'{o}n del rendimiento. \\
\addlinespace
Forward-fill & Propagaci\'{o}n del \'{u}ltimo valor v\'{a}lido hacia delante en el tiempo para rellenar valores ausentes. Preserva la causalidad. \\
\addlinespace
F1-score & Media arm\'{o}nica de Precision y Recall: $F1 = 2PR/(P+R)$. M\'{e}trica principal para clasificaci\'{o}n desbalanceada. \\
\addlinespace
Seq. bias & Sesgo por secuencia: diferencia sistem\'{a}tica en los valores iniciales de las se\~{n}ales entre experimentos distintos. Se elimina restando $x_0$. \\
\addlinespace
Bootstrap & Submuestreo con reemplazamiento: cada \'{a}rbol de RF se entrena en una muestra bootstrap ($\approx 63\%$ \'{u}nico) del training. \\
\addlinespace
Elastic-net & Regularizaci\'{o}n combinada L1 ($\ell_1$-norm) + L2 ($\ell_2$-norm): combina sparsidad (L1) y estabilidad (L2). \\
\bottomrule
\end{tabularx}
\end{center}

\vspace{1cm}
\noindent\textit{Este documento ha sido elaborado con an\'{a}lisis profundos del ``por qu\'{e}'' de cada
decisi\'{o}n t\'{e}cnica, a partir de los resultados del notebook de TD6 del Data Challenge de
Kaggle sobre detecci\'{o}n de fallos en motores servo-rob\'{o}ticos.}

\end{document}
